\chapter{Princípio de funcionamento, ciclos e estados de operação}
\label{cap:funcionamento}

\begin{objetivos}
Ao final deste capítulo você deve ser capaz de:
\begin{itemize}
  \item descrever a rotina de inicialização executada antes do primeiro ciclo;
  \item descrever as três etapas do ciclo de varredura;
  \item explicar o que é tempo de varredura e quais fatores o afetam;
  \item justificar o \emph{watchdog} e reconhecer a falha que ele sinaliza;
  \item diferenciar processamento cíclico, por interrupção, por tempo e por evento;
  \item interpretar o estado dos pontos nas imagens de entrada e de saída;
  \item distinguir os estados de operação do controlador.
\end{itemize}
\end{objetivos}

\section{Rotina de inicialização}

Antes de executar o programa do usuário, o controlador cumpre uma rotina gravada em seu sistema
operacional. Ao ser energizado no estado de execução, ele:

\begin{enumerate}
  \item limpa a memória imagem;
  \item testa a memória RAM;
  \item verifica a possibilidade de executar o programa gravado.
\end{enumerate}

Só depois disso a CPU passa a varrer as instruções em laço, indefinidamente --- é o ciclo de
varredura descrito na seção seguinte, repetido enquanto o controlador estiver em execução.

\section{O ciclo de varredura}

O controlador funciona como um microcomputador dedicado, executando continuamente ciclos de varredura.
O ciclo se divide em três etapas:

\begin{figure}[H]
  \centering
  \includegraphics[width=0.86\linewidth]{09-ciclo-varredura.pdf}
  \caption{As três etapas do ciclo de varredura e a supervisão de tempo.}
  \label{fig:ciclo-varredura}
\end{figure}

\begin{enumerate}
  \item \textbf{Entrada} --- os dados dos sensores e transdutores são lidos e armazenados na
        \emph{memória imagem das entradas};
  \item \textbf{Processamento} --- o programa do usuário é executado sobre esse conteúdo;
  \item \textbf{Saída} --- os resultados são gravados na \emph{memória imagem das saídas} e depois
        transferidos aos cartões de saída, que acionam os atuadores.
\end{enumerate}

O ciclo completo é chamado de \emph{scan}, varredura ou \emph{loop}. O tempo para percorrê-lo é o
\textbf{tempo de varredura} (ou tempo de \emph{scan}), e depende de dois fatores: o tamanho do
programa do usuário e a quantidade de pontos de entrada e saída. Valores típicos são da ordem de
poucos milissegundos; com a eletrônica embarcada atual, controladores rápidos chegam a dezenas de
microssegundos para programas pequenos.

\index{ciclo de varredura}\index{tempo de varredura}\index{memória imagem}

\section{Por que o tempo de varredura importa}

Se o tempo de varredura é maior que a duração de um pulso de entrada, o pulso pode ser perdido: a
entrada vai de 0 para 1 e volta a 0 enquanto o controlador ainda está processando, e nenhum ciclo
chega a ler esse estado. Se a informação de entrada se perde, a saída correspondente não é gerada.

\paragraph{Exemplo numérico.} Um sensor de fim de curso gera um pulso de 5\,ms na passagem da peça.
Com tempo de varredura de 20\,ms, existe chance real de o pulso cair entre duas leituras e
simplesmente não ser visto. Duas soluções possíveis: aumentar a velocidade de processamento
(controlador mais rápido ou programa mais enxuto) ou usar cartão de \emph{contagem rápida} ou entrada
com interrupção, que captura o evento fora do ciclo --- o assunto da Seção seguinte.

Por isso o tempo de varredura é critério de escolha do controlador (Capítulo~\ref{cap:porte}) e
critério de revisão do programa (Capítulo~\ref{cap:boas-praticas}).

\subsection{Watchdog}

Para evitar comportamento indefinido quando o processamento se estende além do previsto, existe um
temporizador de supervisão, o \emph{watchdog time}. Ele supervisiona o tempo de processamento;
ultrapassado o valor máximo --- tipicamente algo em torno de 150\,ms, mas definido por fabricante ---
o funcionamento da CPU é interrompido imediatamente, o controlador assume estado de erro e sinaliza a
falha de forma visual e, às vezes, sonora.

\begin{atencao}
Um controle que reinicia no meio da produção não é um defeito misterioso: é frequentemente watchdog
disparado. As causas mais comuns são laço de programa mal construído, temporização muito longa e
programa grande demais para o tempo configurado. A primeira providência é ler o diagnóstico do
controlador, não reiniciar o equipamento até o problema voltar.
\end{atencao}

\section{Modos de processamento}

Além do modo cíclico, existem três modos que atendem situações em que o ciclo normal não é adequado:

\begin{table}[H]
  \centering
  \caption{Modos de processamento do programa.}
  \label{tab:modos-processamento}
  \begin{tabular}{p{3.2cm} p{6.2cm} p{5.2cm}}
    \toprule
    Modo & Como funciona & Quando é necessário \\
    \midrule
    Cíclico
      & As instruções são lidas sequencialmente do início ao fim e o ciclo recomeça
      & Modo predominante em praticamente todas as CPUs \\
    Por interrupção
      & Ao reconhecer uma ocorrência prioritária, a CPU interrompe o ciclo e executa uma rotina de
interrupção; ao terminar, retoma o ponto onde parou
      & Eventos que não podem esperar o ciclo completo, como emergência ou captura de pulso \\
    Comandado por tempo
      & Trechos do programa são executados em intervalos regulares de tempo, dentro do ciclo normal
      & Controle com amostragem periódica fixa, na ordem de milissegundos \\
    Por evento
      & O programa reage a eventos específicos: retorno de energia, falha de bateria em controlador
legado e estouro do tempo de supervisão do ciclo
      & Procedimentos de inicialização, diagnósticos e resposta a falha \\
    \bottomrule
  \end{tabular}
\end{table}

A distinção tem consequência prática: lógica de segurança ou de captura de evento rápido não deve ser
resolvida no modo cíclico comum.

\section{As imagens de entrada e de saída}

Durante o ciclo, a comunicação com o campo acontece pela memória imagem. Um exemplo mostra por que
isso é importante para ler o programa de outra pessoa.

Considere um trecho de Ladder com um contato NA da entrada \texttt{IN00} em série com um contato NF da
entrada \texttt{IN02}, acionando a saída \texttt{OUT05}. No instante analisado, \texttt{IN00} está
acionada (nível lógico 1) e \texttt{IN02} está desativada (nível lógico 0).

\begin{table}[H]
  \centering
  \caption{Estado dos pontos no instante analisado e resultado do processamento.}
  \label{tab:imagem-es}
  \begin{tabular}{l c c l}
    \toprule
    Ponto & Estado de campo & Bit na imagem & Efeito no programa \\
    \midrule
    IN00 & contato fechado & 1 & contato NA conduz \\
    IN01 & contato aberto  & 0 & --- \\
    IN02 & contato aberto  & 0 & contato NF conduz (lógica de negação) \\
    \midrule
    OUT05 & --- & 1 & saída acionada \\
    \bottomrule
  \end{tabular}
\end{table}

Em lógica combinacional, o que o programa implementa é:

$$
\texttt{OUT05} = \texttt{IN00} \cdot \overline{\texttt{IN02}}
$$

A saída é acionada quando o nível das \emph{imagens} obedece a essa relação. Note o ponto que costuma
confundir: o bit da imagem depende de o ponto estar energizado, e não de o contato de campo ser NA ou
NF. Um contato NF de campo energizado resulta em bit 1; a inversão que se vê no programa vem do tipo
de contato usado na lógica, não do campo.

\begin{pratica}
Usando um simulador de CLP ou a função de simulação do \emph{software} de programação:
(a) monte o trecho do exemplo e force \texttt{IN00} e \texttt{IN02} para reproduzir a
Tabela~\ref{tab:imagem-es};
(b) monte uma tabela-verdade completa com as quatro combinações de \texttt{IN00} e \texttt{IN02} e
confira com a expressão;
(c) troque o contato NF por um NA e explique, por escrito, em que situação a saída passa a se comportar
de forma perigosa.
\end{pratica}

\section{Estados de operação}

A UCP pode assumir um estado de erro e dois estados de programação:

\begin{figure}[H]
  \centering
  \includegraphics[width=\linewidth]{09-estados.pdf}
  \caption{Estados de operação do controlador e transições.}
  \label{fig:estados}
\end{figure}

\paragraph{Programação (STOP).} O controlador não executa o programa e não assume nenhuma lógica de
controle. Fica pronto para ser configurado, receber programa novo ou ser modificado. A programação
feita nesse estado é chamada \emph{off-line}.

\paragraph{Execução (RUN).} O controlador executa o programa do usuário. Em alguns controladores é
possível alterar o programa com ele em execução --- programação \emph{on-line}. Quando a lógica
modificada é enviada, o novo programa entra em operação sem que o estado mude e sem interromper o
ciclo.

\paragraph{Erro.} Estado que aponta falha de operação ou de execução do programa.

\subsection{Auto-teste e proteção de dados}

A cada inicialização de ciclo o controlador executa um auto-teste do circuito interno; se houver
falha --- memória ou microprocessador, por exemplo --- um sinal de alarme é emitido. Em caso de falta
de energia, as memórias são alimentadas automaticamente pelo sistema por baterias de reserva, para
evitar perda de dados.

\begin{atencao}
Programação \emph{on-line} é recurso poderoso e perigoso. Alterar lógica com a máquina produzindo
exige: autorização formal, procedimento de bloqueio quando o trecho for de segurança, e registro da
alteração feita. Sem isso, o próximo mantenedor não saberá qual programa está efetivamente na
memória --- exatamente o problema que o Capítulo~\ref{cap:documentacao} trata.
\end{atencao}

\begin{exercicios}
\begin{enumerate}[label=\alph*)]
  \item Descreva as três etapas do ciclo de varredura indicando, em cada uma, onde o dado está
        armazenado.
  \item Um programa tem tempo de varredura de 30\,ms e precisa detectar um pulso de 8\,ms. Que
        soluções existem, e qual delas não exige trocar o controlador?
  \item Explique a função do \emph{watchdog} e o que o operador deve fazer quando ele dispara.
  \item Por que o watchdog é classificado como processamento por evento?
  \item Calcule o valor de \texttt{OUT05} para as quatro combinações de \texttt{IN00} e
        \texttt{IN02} e monte a tabela-verdade.
  \item Um técnico afirma: ``o contato NF de campo apareceu como bit 1, então o programa está
        errado''. Corrija a afirmação.
  \item Explique a diferença entre programação \emph{off-line} e \emph{on-line}, e uma consequência
        de controle de qualidade de cada uma.
  \item Por que o auto-teste de RAM acontece antes do ciclo, e não durante?
\end{enumerate}
\end{exercicios}
